\section{The Church that was Left Behind}

For many Christians, referring to a ``secret Gospel'' is either presumptuous, dangerous, or absurd --- perhaps a bit of all three. And, indeed, such fears are usually the case with those who go in quest of Gnostic hierophanies\footnote{\url{https://en.wikipedia.org/wiki/Hierophany}} -- just witness the popular mesmerization over the Da Vinci code novels and the ``secret Gospels'' (as if someone could read them \& suddenly become enlightened with an alternative dogma-system). However, as Gornahoor has endeavored to demonstrate\footnote{\url{http://www.gornahoor.net/?p=2091\#more-2091}}, there are important episodes of early Church history which intimate that much was lost (as perhaps, gained) during the \emph{Ekklesia}`s struggle to acrete and purify the core dogmas. Clement's \emph{Stromata} represent such an effort:

\begin{quotationx}
Again he says that those who are ``still blind and dumb, not having understanding, or the un-dazzled and keen vision of the contemplative soul . . . must stand outside of the divine choir. . . . Wherefore, in accordance with the method of concealment, the truly sacred Word, truly divine and most necessary for us, deposited in the shrine of truth, was by the Egyptians indicated by what were called among them adyta, and by the Hebrews by the veil. Only the consecrated . . . were allowed access to them. For Plato also thought it not lawful for ` the impure to touch the pure. Thence the prophecies and oracles are spoken in enigmas, and the Mysteries are not exhibited incontinently to all and sundry, but only after certain Page 69 purifications and previous instructions''. Ibid., bk.V, ch.iv. He then descants at great length on Symbols, expounding Pythagorean, Hebrew, Egyptian, Ibid, ch. v-viii and then remarks that the ignorant and unlearned man fails in understanding them. ``But the Gnostic apprehends. Now then it is not wished that all things should be exposed indiscriminately to all and sundry, or the benefits of wisdom communicated to those who have not even in a dream been purified in soul (for it is not allowed to hand to every chance comer what has been procured with such laborious efforts); nor are the Mysteries of the Word to be expounded to the profane''. The Pythagoreans and Plato, Zeno, and Aristotle had exoteric and esoteric teachings. The philosophers established the Mysteries, for ``was it not more beneficial for the holy and blessed contemplation of realities to be concealed?'' Ibid., ch. ix. The Apostles also approved of ``veiling the Mysteries of the Faith'', ``for there is an instruction to the perfect'', alluded to in Colossians i, 9-11 and 25-27.

So that, on the one hand, then, there are the Mysteries which were hid till the time of the Apostles, and were delivered by them as they were received from the Lord, Page 70 and, concealed in the Old Testament, were manifested to the saints. And, on the other hand, there is ` the riches of the glory of the mystery in the Gentiles,' which is faith and hope in Christ; which in another place he has called the ``foundation''. He quotes S. Paul to show that this ``knowledge belongs not to all'', and says, referring to Heb. v. and vi., that ``\emph{there were certainly among the Hebrews, some things delivered unwritten}''; and then refers to S. Barnabas, who speaks of God, ``who has put into our hearts wisdom and the understanding of His secrets'', and says that ``it is but for few to comprehend these things'', as showing a ``trace of Gnostic tradition''.''Wherefore instruction, which reveals hidden things, is called illumination, as it is the teacher only who uncovers the lid of the ark''. Ibid., bk. V, ch. x Further referring to S. Paul, he comments on his remark to the Romans that he will ``come in the fullness of the blessing of Christ, Loc. Cit, XX, 29. '' a and says that he thus designates ``the spiritual gift and the Gnostic interpretation, which being present he desires to impart to them present as ` the fullness of Christ, according to the revelation of the Mystery sealed in the ages of eternity, but now manifested by the prophetic Page 71 Scriptures' Ibid., xvi;ten, 25-26; the version quoted differs in words, but not in meaning, from the English Authorised Version ... ..But only to a few of them is shown what those things are which are contained in the Mystery. Rightly, then, Plato, in the epistles, treating of God, says: ` \emph{We must speak in enigmas; that should the tablet come by any mischance on its leaves either by sea or land,he who reads may remain ignorant}`.'' Stromata, bk. V, ch. x After much examination of Greek writers, and an investigation into philosophy, S. Clement declares that the Gnosis ``imparted and revealed by the Son of God, is wisdom. . . . And the Gnosis itself is that which has descended by transmission to a few, having been imparted unwritten by the Apostles''. Ibid., bk. VI, ch. vii A very long exposition of the life of the Gnostic, the Initiate, is given, and S. Clement concludes it by saying: ``Let the specimen suffice to those who have ears. For it is not required to unfold the mystery, but only to indicate what is sufficient for those who are partakers in knowledge to bring it to mind''. Ibid., bk. VII, ch. xiv. Regarding Scripture as consisting of allegories and symbols, and as hiding the sense in order to Page 72 stimulate enquiry and to preserve the ignorant from danger. Ibid., bk. VI, ch. xv. S. Clement naturally confined the higher instruction to the learned. ``Our Gnostic will be deeply learned'', Ibid., bk. VI, x he says. ``Now the Gnostic must be erudite''. Ibid., bk. VI, vii Those who had acquired readiness by previous training could master the deeper knowledge, for though ``a man can be a believer without learning, so also we assert that it is impossible for a man without learning to comprehend the things which are declared in the faith''. Ibid., bk. I,ch. vi ``Some who think themselves naturally gifted, do not wish to touch either philosophy or logic; nay more, they do not wish to learn natural science. They demand bare faith alone. . . So also I call him truly learned who brings everything to bear on the truth --- so that, from geometry, and music, and grammar, and philosophy itself, culling what is useful, he guards the faith against assault. How necessary is it for him who desires to be partaker of the power of God, to treat of intellectual subjects by philosophising''. Ibid., ch. ix. ''The Gnostic avails himself of branches of learning as auxiliary Page 73 preparatory exercise.'' Ibid., BK. VI, ch. x. So far was S. Clement from thinking that the teaching of Christianity should be measured by the ignorance of the unlearned. ``He who is conversant with all kinds of wisdom will be pre-eminently a Gnostic''. Ibid., bk. I, ch. xiii. '' Thus while he welcomed the ignorant and the sinner, and found in the Gospel what was suited to their needs, he considered that only the learned and the pure were fit candidates for the Mysteries. ``The Apostle, in contradistinction to Gnostic perfection, calls the common faith the foundation, and sometimes milk'', Vol.XII. Stromata, bk. V, ch. iv. but on that foundation the edifice of the Gnosis was to be raised, and the food of men was to succeed that of babes.
\end{quotationx}

True Gnosticism was (in other words) an entrance to the ``Greater Mysteries'' which either could (or perhaps ``had'') to begin in the ``Lesser Mysteries''; it was to be accomplished by an elite of spirit \& mind, and would involve the re-integration of all human learning into Hermetic (ie., underground) frameworks.

These are the commentaries of Annie Besant\footnote{\url{http://www.theosophical.ca/books/EsotericChristianity\_AnnieBesant.pdf}}, a theosophical\footnote{\url{https://en.wikipedia.org/wiki/Annie\_Besant}} writer from the last century. It is discouraging to see that such things have been brought to the attention of the educated Christian public long before, but apparently with only the result that confidence in the Faith has been undermined, without putting anything in its place (presumably this is why Evola disapproved of Steiner \& Besant). Gornahoor has emphasized (with Evola) that it is ``not a matter of showing the same misunderstanding towards Christianity as Christianity has often showed to everything else''. When A. Besant was struggling to keep her faith \& requested more books to assist her, the leader\footnote{\url{https://en.wikipedia.org/wiki/Edward\_Bouverie\_Pusey}} of the Catholic Anglican Revival told her ``she had read too much already''.

Likewise, the same obtuseness is encountered from Eastern Orthodox Christians when Origen, Clement, or John Philoponus is brought up -- of all the Big Three branches of the Faith, one would think the Orthodox would be sympathetic to those who are given ``bad press'', even if by a Church Council. However, this is how it stands -- the Church will not budge when one begins to read the Bible like the Bereans, and inquires as to the ``mysteries'' of the Faith -- they are told that they are perfectly embodied (and imperfectly understood) in the completed Dogmas. Apparently, there are even accounts available which trace back to Peter's\footnote{\url{http://jacksonsnyder.com/arc/bible-news/apostolia/pdf/ALL-CLEMENT-6x9-120707-embedded.pdf}} journeys (although they cannot possibly stem from Clement himself, who was born many years after). Given the resurfacing of the Nag Hammadi hoard (immediately after the apocalyptic World Wars), as well as existence of such monuments as the \emph{Stromata}, the Church's behavior is utterly inexplicable. To sum it up perfectly, the Church has declared \emph{by fiat} that exoteric truth is the only kind of truth which exists.

It seems intent on preserving an orthodoxy which has been shown (to even the most casual eye) to have been largely a product of intense control, a control arguably necessary both then and even now, but which cannot possibly be exhaustive. The only possible route which could save the Ekklesia is to explore the very foundations of the Faith ``once delivered to the saints'' and see if (possibly) the ``living ground and pillar of the Truth'' could refer, not necessarily to those who write the official Church histories, but to those who have actually helped co-create or sub-create a Reality which is responsive to the truly transcendent and divine image in man, an image both solar and masculine reached for in the Book of Revelation\footnote{\url{https://www.youtube.com/watch?v=k9IfHDi-2EA}}, written by the apostle of the eagle\footnote{\url{https://www.amazon.com/Voice-Eagle-Heart-Celtic-Christianity/dp/0970109709}} -- John.

It is only in the light of such a Sun that the lesser lights and mysteries (such as the Moon) can have their own lumination. Does the \emph{Ekklesia} wish to remain forever a vehicle at best of the lesser mysteries, or does it wish to take its rightful place as the inheritor of the secret fire\footnote{\url{https://en.wikipedia.org/wiki/Sacred\_fire\_of\_Vesta}}? If it does not, its own Scripture testifies against it, that it will be spit out as lukewarm, or even caught up in the heresy of the Anti-Christ\footnote{\url{http://www.berdyaev.com/berdiaev/berd\_lib/1913\_170.html}} (a fear which dominated the thoughts of the long-suffering Russian Church of the last century) --

\begin{quotex}
For the mystery of iniquity doth already work: only he who now letteth will let, until he be taken out of the way.
\end{quotex}


Surely, if there was an earthly Jerusalem, and she is fallen, we have a greater -- a heavenly Jerusalem, the true mother above. \emph{The Protestant whim about the ``invisible Church'' may be more true than they imagine}. And if the chosen ones called Jews were cut off to graft in the Gentiles, how much more should the \emph{Ekklesia} be jealous to guard the true sanctity of her mysteries, lest the Church itself be slain, since judgement begins with the house of God? Rather than attempting to define itself against a tradition which distanced itself from Divine Tradition as it unfolded (the Jews), the modern Church has pursued a philo-Semitic course; in another context (that of real Judaic tradition) this might be worthwhile, yet it places the \emph{Ekklesia} in the position of sharing the same literal, awkward interpretation of the Scriptures which men like Origen\footnote{\url{https://books.google.com/books?id=UC0eWGKHa9YC&pg=PA112\&lpg=PA112\&dq=CLEMENT+SECRET+GNOSIS\&source=bl\&ots=PKa9DTVJav\&sig=ODd0rSifdZ2VRT-MLg\_LMI3BInA\&hl=en\&ei=i6HZTZb2C8nq0QG25Jz8Aw\&sa=X\&oi=book\_result\&ct=result\&resnum=3\&ved=0CBwQ6AEwAjgK\#v=onepage\&q=origen\&f=false}} attributed to the Semitic spirit, and which the early Fathers ridiculed in Jew and pagan alike.


\flrightit{Posted on 2011-09-02 by Logres}

